\begin{abstract}
We introduce \textsc{QLoRA}, a memory-efficient finetuning technique that enables the finetuning of large-scale models with up to 65 billion parameters on a single 48GB GPU, all while maintaining the performance standards associated with full 16-bit finetuning tasks. \textsc{QLoRA} operates by propagating gradients through a fixed, 4-bit quantized version of a pretrained language model, directing updates into Low Rank Adapters~(LoRA). Our leading model suite, termed \textbf{Guanaco}, surpasses all previously available open models on the Vicuna benchmark, achieving 99.3\% of ChatGPT’s performance with just one day of finetuning on a single GPU. The innovations underpinning \textsc{QLoRA} include: (a) 4-bit NormalFloat (NF4), a new data format tailored for optimal information retention under normally distributed weights, (b) Double Quantization, which further lowers memory usage by applying quantization to quantization parameters, and (c) Paged Optimizers, designed to smooth out memory consumption spikes. Leveraging \textsc{QLoRA}, we successfully finetune over 1,000 models and present comprehensive analyses on instruction-following and chatbot capabilities, spanning 8 instruction datasets, various architectures (LLaMA, T5), and scales beyond the practical reach of standard finetuning approaches (such as 33B and 65B parameter configurations). Our experiments reveal that QLoRA-based finetuning on concise, high-quality datasets yields state-of-the-art outcomes, often outperforming prior leading models even at reduced model sizes. We present an in-depth evaluation of chatbot responses using both human and GPT-4 judgments, highlighting that GPT-4 provides a reliable and cost-effective substitute for human raters. Moreover, our findings suggest that current chatbot assessment metrics do not provide reliable insights into actual chatbot proficiency. Through targeted analyses, we also identify scenarios where \textbf{Guanaco} lags behind ChatGPT. All model checkpoints and implementation code are publicly available, along with custom CUDA kernels supporting 4-bit training.\footnote{\url{https://github.com/artidoro/qlora} and \url{https://github.com/TimDettmers/bitsandbytes}}
\end{abstract}

\section{Introduction}

Adapting large language models (LLMs) through finetuning has proven to be a particularly powerful approach for enhancing model capabilities~\citep{min2021metaicl, wei2021finetuned, ouyang2022training, wang2022super, wang2022self, liu2022few} and for tailoring behaviors, whether by reinforcing favorable attributes or suppressing unwanted ones~\citep{ouyang2022training,askell2021general,bai2022training}. Nonetheless, finetuning models at massive scales remains largely inaccessible due to the significant computational and memory costs involved. For instance, conventional 16-bit finetuning of a LLaMA model with 65 billion parameters~\citep{touvron2023llama} requires in excess of 780 GB of GPU RAM. Quantization strategies have emerged to decrease the memory required for LLM inference~\citep{dettmers2022llm, dettmers2022case, frantar2022gptq, xiao2022smoothquant}; however, these approaches have so far not extended to training, as they frequently encounter instability during that phase~\citep{wortsman2023stable}.

In this work, we establish for the first time that it is feasible to finetune a model quantized to 4 bits without any drop in model efficacy. Our proposed method, \textsc{QLoRA}, employs an innovative high-precision quantization mechanism for compressing a pretrained model to 4 bits, subsequently augmenting the model with a compact set of trainable Low-rank Adapter parameters~\citep{hu2021lora} 

which are updated through gradient backpropagation across the quantized model weights. 

\textsc{QLoRA} brings down the average GPU memory needed for finetuning a 65B parameter model from over 780GB to under 48GB, with no negative impact on speed or prediction quality compared to full 16-bit finetuning. This represents a substantial advancement in democratizing LLM finetuning: even the largest open models can now be adapted on a single GPU. Leveraging \textsc{QLoRA}, we introduce the \textbf{Guanaco} series of models, with the second largest achieving 97.8\% of ChatGPT’s performance on the Vicuna~\citep{vicuna2023} evaluation, requiring less than 12 hours of training on a standard consumer GPU; the largest model attains 99.3\% within 24 hours on a single high-end professional GPU, thus closing the performance gap with ChatGPT on the Vicuna benchmark. For inference, our smallest \textbf{Guanaco} variant (7B parameters) only uses 5 GB of memory and surpasses the 26 GB Alpaca model by more than 20 points on the Vicuna benchmark (see Table~\ref{tbl:vicuna_eval}).

\textsc{QLoRA} presents a set of advancements targeting reduced memory consumption while preserving model accuracy: (1) \textbf{4-bit NormalFloat}, a quantization format specifically optimized for normally distributed data from an information-theoretic perspective, achieves superior empirical performance relative to both 4-bit Integer and 4-bit Float representations. (2) \textbf{Double Quantization}, a technique in which the quantization coefficients themselves are quantized, resulting in an average savings of approximately 0.37 bits per parameter—equivalent to nearly 3 GB for models containing 65B parameters. (3) \textbf{Paged Optimizers}, which leverage NVIDIA's unified memory to circumvent the memory spikes associated with gradient checkpointing during the processing of long-sequence mini-batches. These elements are integrated into an enhanced LoRA framework that places adapters at every network layer, thereby mitigating the accuracy compromises observed in previous studies.

The resource efficiency of \textsc{QLoRA} allows for a comprehensive investigation of instruction finetuning and chatbot efficacy on model scales that would be impractical with conventional finetuning methods due to excessive memory demands. Consequently, we trained over 1,000 models spanning multiple instruction tuning datasets, architectural configurations, and ranging in size from 80M to 65B parameters. Our findings demonstrate not only that \textsc{QLoRA} can match the performance of standard 16-bit training (\S\ref{sec:comp_to_std}) and facilitate the training of state-of-the-art systems such as the \textbf{Guanaco} chatbot (\S\ref{sec:pushing}), but also reveal patterns across trained models. Notably, our results indicate that the quality of training data surpasses sheer dataset scale in importance: for instance, a 9,000-sample dataset (OASST1) led to better chatbot outcomes than a 450,000-sample dataset (FLAN v2, subsampled), even when both aimed to promote generalizable instruction following. Moreover, we establish that high performance on the Massive Multitask Language Understanding (MMLU) benchmark does not guarantee strong results on the Vicuna chatbot benchmark, and vice versa—highlighting that the appropriateness of the dataset to the target application is more critical than its size.

In addition, we conduct a comprehensive evaluation of chatbot systems that involves assessments from both human judges and GPT-4. Our methodology employs a tournament-based benchmarking approach, where models are paired to generate responses to the same prompt and are then evaluated against each other. Each round’s winner is decided by either human annotators or GPT-4. The aggregated outcomes from these matchups are synthesized into Elo ratings~\citep{elo1967proposed, elo1978rating}, thereby establishing the hierarchy of model capabilities. We observe substantial alignment between human judgments and GPT-4 evaluations regarding overall model rankings; however, notable discrepancies are also present in some cases. This underscores the fact that, despite being more cost-effective, model-based assessments introduce additional sources of uncertainty compared to human annotation.

To further contextualize our benchmark findings, we provide an in-depth qualitative review of the \textbf{Guanaco} models, drawing attention to particular strengths and weaknesses not reflected by quantitative scores.

For transparency and to promote further research, we make publicly available all chatbot outputs accompanied by both human and GPT-4 annotations. Our complete codebase, including CUDA kernels, is released as open source, and we incorporate our tools into the Hugging Face transformers framework~\citep{wolf2019huggingface}, simplifying their adoption. Additionally, we provide adapter collections for 7/13/33/65B model sizes, each trained on eight different instruction tuning datasets, resulting in a total of 32 unique fine-tuned, open-source models.

\section{Background}
\paragraph{Block-wise k-bit Quantization} Quantization refers to the procedure of mapping values from a high-precision representation to one with reduced precision. Frequently, this involves converting from a data type with a greater bit-width—such as 32-bit floating point numbers—to one with fewer bits, such as 8-bit integers. To maximize the utilization of the representational capacity in the lower-bit format, the original data are typically normalized relative to their maximum absolute value, generally computed across tensor elements. As an illustration, when converting a FP32 tensor to an Int8 tensor spanning $[-127, 127]$, the transformation is:

\begin{equation}
    \mathbf{X}^{ \text{Int8}} = \text{round}\left( \frac{127}{{\text{absmax}}(\mathbf{X}^{{\text{FP32}}})}\mathbf{X}^{{\text{FP32}}} \right) = \text{round}(c^{\text{FP32}}\cdot \mathbf{X}^{{\text{FP32}}}),
\end{equation}

with $c$ denoting the {\em quantization constant} or {\em scale factor}. The corresponding reverse process, dequantization, is defined as:

\begin{equation}
    \text{dequant}(c^{\text{FP32}}, \mathbf{X}^{\text{Int8}}) = \frac{\mathbf{X}^{\text{Int8}}}{c^{\text{FP32}}} = \mathbf{X}^{\text{FP32}}
\end{equation}

A significant drawback of this method arises when the input tensor contains a value of high magnitude (an outlier); in such cases, the quantization bins—specific combinations of bits—tend to be underutilized, with some bins containing few or no quantized values. To address this issue, a widely-used strategy is to partition the input tensor into smaller segments, or blocks, each of which is quantized separately and assigned its own quantization constant $c$. This can be expressed more precisely as follows: The input tensor $\mathbf{X}\in \mathbb{R}^{b\times h}$ is divided into $n$ consecutive blocks of size $B$ by flattening the tensor and slicing the resulting array into ${n = ({b\times h})/{B}}$ blocks. Each of these blocks is then quantized in isolation using Equation~1, resulting in a quantized tensor accompanied by $n$ block-specific quantization constants $c_i$.

\paragraph{Low-rank Adapters} The Low-rank Adapter (LoRA) finetuning technique \citep{hu2021lora} offers a means of reducing memory consumption by introducing a compact set of trainable adapter parameters, while the original model weights remain frozen. During stochastic gradient descent, gradients propagate through the unmodified pretrained model parameters to the adapters, which are updated in accordance with the loss objective. LoRA enhances a standard linear projection by incorporating an additional low-rank decomposition. Consider a projection of the form ${\mathbf{X} \mathbf{W} = \mathbf{Y}}$, where $\mathbf{X}\in  \mathbb{R}^{b\times h}$ and $\mathbf{W}\in  \mathbb{R}^{h\times o}$; in LoRA, the operation becomes:

\begin{equation}
    \mathbf{Y} = \mathbf{X}\mathbf{W} + s\mathbf{X}\mathbf{L}_1\mathbf{L}_2,
\end{equation}

where $\mathbf{L}_1\in  \mathbb{R}^{h\times r}$ and $\mathbf{L}_2\in  \mathbb{R}^{r\times o}$, and $s$ denotes a scalar value.

\paragraph{Memory Requirement of Parameter-Efficient Finetuning} The discussion around the memory usage of LoRA in training primarily centers on both the quantity and dimensionality of adapters utilized. Due to LoRA's extremely lightweight memory demands, it is feasible to deploy a larger number of adapters to boost model performance while barely affecting the overall memory consumption. Although LoRA is a Parameter Efficient Finetuning (PEFT) strategy, in the context of finetuning LLMs, the dominant memory load comes from activation gradients rather than the parameters introduced by LoRA itself. For example, when applying LoRA to a 7B LLaMA model finetuned on FLAN v2 with batch size set to 1, where the LoRA parameters correspond to approximately 0.2\% of the model's total parameters as in prior works~\citep{hu2021lora,liu2022few}, the memory needed for storing the input gradients to LoRA is 567 MB, in contrast to only 26 MB required for the LoRA parameters themselves. Utilizing gradient checkpointing~\citep{chen2016training}, the input gradient storage per sequence drops significantly to around 18 MB, which still surpasses the total memory held by all LoRA weights. For further comparison, the memory usage of the base model at 4-bit quantization stands at 5,048 MB. These observations underscore the relevance of gradient checkpointing for memory optimization, while also suggesting that further compression of LoRA parameters yields limited reductions in memory usage. Consequently, increasing the number of adapters remains practical, with minimal effect on the memory demands during training (refer to Appendix~\ref{app:memory} for an in-depth analysis). This flexibility plays a vital role in achieving performance that approaches full 16-bit precision, as will be examined in subsequent sections.

\section{\textsc{QLoRA} Finetuning}

\textsc{QLoRA} facilitates accurate 4-bit finetuning by leveraging two novel strategies: 4-bit NormalFloat (NF4) quantization and Double Quantization. The method also incorporates Paged Optimizers to address the issue of memory surges during gradient checkpointing, a common source of out-of-memory failures that have long hindered the finetuning of large-scale models on a single computing node.

Within the \textsc{QLoRA} framework, one data type is reserved for storage—typically in 4-bit precision—while computations are executed using a higher precision format, generally BFloat16. Concretely, each time a \textsc{QLoRA} model weight tensor is accessed, it undergoes dequantization to BFloat16, after which matrix multiplications are carried out in this 16-bit format.

The subsequent sections detail each component of \textsc{QLoRA} and provide a formal specification.

\paragraph{4-bit NormalFloat Quantization} The NormalFloat (NF) format extends Quantile Quantization \citep{dettmers2022optimizers}, a data type designed for information-theoretic efficiency by ensuring an even assignment of input tensor values to quantization bins. Quantile quantization operates by calculating tensor quantiles via the empirical cumulative distribution function.

A significant drawback of quantile quantization is the computational overhead associated with accurate quantile computation. To mitigate this, efficient quantile estimation techniques—such as SRAM quantiles~\citep{dettmers2022optimizers}—are often utilized. Nevertheless, due to the inherent approximation, these approaches tend to incur substantial quantization errors for outlier values, which are typically critical.

When input tensors are drawn from a distribution that remains invariant up to a scaling constant for quantization, both high computational cost and approximation inaccuracies can be alleviated. Under these circumstances, the identical quantiles across tensors allow for precise quantile computation with tractable effort.

Given that pretrained neural network weights typically follow a zero-centered normal distribution with standard deviation $\sigma$ (see Appendix~\ref{app:normality}), it is possible to unify all weight distributions by appropriately scaling $\sigma$ so that the resultant distribution precisely aligns with the representable range of the selected data type. Here, we designate the fixed interval $[-1, 1]$ as the operational range of our data type. Consequently, it becomes necessary to normalize both the quantiles associated with the data type and those from the neural network weights to this shared interval.

The information-theoretically optimal quantization scheme for zero-mean normal distributions (with varying $\sigma$) mapped to $[-1, 1]$ involves the following: (1) determine the $2^k + 1$ quantile values for the standard normal distribution $N(0, 1)$, thus defining a $k$-bit quantile-based quantization format for normal distributions; (2) scale the resulting data type quantiles so that their values fall within $[-1, 1]$; and (3) apply absolute maximum rescaling to normalize any input weight tensor to this same interval, enabling consistent quantization.

When the scales of the weights and the quantization data type coincide, the quantization process can proceed in the standard manner. Notably, step (3) corresponds to adjusting the standard deviation of the weights to that of the $k$-bit data type. Formally, the $2^k$ representative quantile values $q_i$ of the quantized data type can be derived by:

\begin{equation}
q_i  = \frac{1}{2}\left( Q_X\left(\frac{i}{2^k+1}\right) + Q_X\left(\frac{i+1}{2^k+1}\right)\right),
\end{equation}

where $Q_X(\cdot)$ represents the quantile function corresponding to the standard normal distribution $N(0,1)$. In the context of symmetric k-bit quantization, a notable limitation is the absence of an explicit encoding for zero, which is critical when quantizing padding or other zero entries to maintain exactness. To guarantee that a discrete zero point aligns with the value $0$ and to fully utilize the $2^k$ discrete values available for a k-bit data type, we construct an asymmetric data type by separately computing quantiles $q_i$ for two intervals: $2^{k-1}$ quantiles for the negative half and $2^{k-1}+1$ for the positive half. These two groups of quantiles are then merged, and the duplicate zero arising from both partitions is eliminated. We refer to this data type, which ensures that each quantization interval contains, in expectation, the same number of samples, as the \emph{k-bit NormalFloat} (NFk). This quantizer achieves optimal information representation for data that is centered at zero and follows a normal distribution. The precise quantized values for this data type are provided in Appendix~\ref{app:nf4}.

\paragraph{Double Quantization{}} 
We propose the technique of \emph{Double Quantization} (DQ), where even the quantization constants themselves undergo quantization to yield further reductions in memory usage. Although high accuracy in 4-bit quantization demands a relatively small blocksize~\citep{dettmers2022case}, such configuration leads to notable memory consumption. For instance, applying 32-bit quantization constants to blocks of size 64 when quantizing $\mathbf{W}$ incurs a storage overhead of $32/64 = 0.5$ bits per parameter. Double Quantization addresses this overhead by compressing the quantization constants.

In detail, Double Quantization considers the initial quantization constants $c_2^{\text{FP32}}$ as inputs for an additional quantization layer. This process produces a set of secondary quantized constants $c_2^{\text{FP8}}$ and new quantization constants $c_1^{\text{FP32}}$ for this second quantization. We implement 8-bit floating point quantization with a blocksize of 256 at this stage, observing no loss in performance with 8-bit quantization, consistent with findings from \citet{dettmers2022case}. Since $c_2^{\text{FP32}}$ are inherently non-negative, we subtract their mean prior to quantization, thereby centering their distribution at zero and permitting symmetric quantization. On average, when employing a blocksize of 64, this approach reduces the per-parameter memory cost from $32/64 = 0.5$ bits to $8/64 + 32/(64\cdot256) = 0.127$ bits, amounting to a memory reduction of 0.373 bits per parameter.

\paragraph{Paged Optimizers} leverage the NVIDIA unified memory~\footnote{\tiny \url{https://docs.nvidia.com/cuda/cuda-c-programming-guide}} capability, which manages page-level data transfers automatically between CPU and GPU. This mechanism ensures seamless GPU operation even in situations where GPU memory is insufficient, similar to how operating system paging manages memory between CPU RAM and storage. By exploiting this feature, our approach allocates optimizer state memory as paged memory, allowing states to be transparently offloaded to CPU RAM when GPU memory is constrained, and then reloaded into GPU memory when optimizer updates require it.

\paragraph{\textsc{QLoRA}.} Building upon the aforementioned components, we specify \textsc{QLoRA} for an individual linear transformation in the quantized backbone model, accompanied by a single LoRA adapter, as shown below:

\begin{equation}
    \mathbf{Y}^{\text{BF16}} =  \mathbf{X}^{\text{BF16}} \text{doubleDequant}(c_1^{\text{FP32}}, c_2^{\text{k-bit}}, \mathbf{W}^{\text{NF4}}) + \mathbf{X}^{\text{BF16}}\mathbf{L}^{\text{BF16}}_1\mathbf{L}^{\text{BF16}}_2,
\end{equation}

where the doubleDequant$(\cdot)$ operation is given by:

\begin{equation}
    \text{doubleDequant}(c_1^{\text{FP32}}, c_2^{\text{k-bit}}, \mathbf{W}^{\text{k-bit}}) = \text{dequant}(\text{dequant}(c_1^{\text{FP32}}, c_2^{\text{k-bit}}), \mathbf{W}^{\text{4bit}}) = \mathbf{W}^{\text{BF16}},
\end{equation}

The quantization format used for $\mathbf{W}$ is NF4, while $c_2$ is stored in FP8.

A block size of 64 is selected for $\mathbf{W}$ to improve quantization fidelity, whereas for $c_2$, a block size of 256 is chosen to optimize memory usage.

During training, parameter updates require computing gradients exclusively with respect to the adapter weights, $\frac{\partial E}{\partial \mathbf{L}_i}$, omitting the calculation of gradients for the 4-bit weight parameters, $\frac{\partial E}{\partial \mathbf{W}}$. Nevertheless, obtaining $\frac{\partial E}{\partial \mathbf{L}_i}$ necessitates evaluating $\frac{\partial \mathbf{X}}{\partial \mathbf{W}}$, which is achieved according to equation~(5) by converting the stored quantized weights $\mathbf{W}^{\text{NF4}}$ to the computation type $\mathbf{W}^{\text{BF16}}$ for accurate derivative computation in BFloat16 format.

In summary, \textsc{QLoRA} utilizes a single storage format—typically 4-bit NormalFloat—and a separate computation format—16-bit BrainFloat. During both forward and backward propagation, the model dequantizes parameters from the storage type to the computation type. However, gradient updates are restricted to the LoRA parameters, which are represented in 16-bit BrainFloat.

\section{QLoRA vs. Standard Finetuning}
\label{sec:comp_to_std}

Previously, we detailed the operational aspects of QLoRA and its ability to drastically minimize memory consumption during model finetuning. The central concern addressed here is whether QLoRA achieves performance comparable to traditional full-model finetuning. Additionally, we explore QLoRA's architectural components, including assessing the effect of NormalFloat4 relative to conventional Float4. The subsequent sections outline experiments designed to investigate these topics.

\paragraph{Experimental setup.} Our study involves three model types—encoder, encoder-decoder, and decoder-only—and systematically compares QLoRA to both 16-bit adapter-finetuning and full-model finetuning for model sizes up to 3B parameters. We evaluate performance on benchmarks including GLUE \citep{wang2018glue} using RoBERTa-large \citep{liu2019roberta}; Super-NaturalInstructions (TKInstruct) \citep{wang2022super} with T5 \citep{t5}; and 5-shot MMLU \citep{hendrycksmeasuring} after finetuning LLaMA models with Flan v2 \citep{longpre2023flan} and Alpaca \citep{alpaca}. To further evaluate the benefits of NF4 over alternative 4-bit formats, we replicate the methodology of \citet{dettmers2022case}, examining post-quantization zero-shot accuracy and perplexity across multiple model families (OPT~\citep{zhang2022opt}, LLaMA~\citep{touvron2023llama}, BLOOM~\citep{scao2022bloom}, and Pythia~\citep{biderman2023pythia}) over a parameter range of 125m to 13B.

Detailed configurations for each experiment are presented in the results section for clarity, with comprehensive experimental information provided in Appendix~\ref{app:exp-setup}.

While paged optimizers are indispensable for enabling 33B/65B \textsc{QLoRA} finetuning on hardware with a single 24/48GB GPU, we omit direct benchmarks for Paged Optimizers because page transfers generally occur only with mini-batches of exceptionally long sequence lengths, which are infrequent. Nonetheless, we assess the runtime for paged optimizers on 65B-parameter models running on 48GB GPUs and observe that, using a batch size of 16, training speeds match those obtained with standard optimizers. We recommend future investigations aim to systematically identify and quantify the specific scenarios in which paging-induced slowdowns manifest.

\paragraph{Default LoRA hyperparameters do not match 16-bit performance} 

Employing conventional LoRA strategies—namely, applying LoRA modules exclusively to the query and value projection matrices of the attention mechanism \citep{hu2021lora}—fails to recover full-model finetuning performance, particularly for larger model backbones. Figure~\ref{fig:hyper}, which displays the results of finetuning LLaMA 7B on Alpaca, indicates that the paramount LoRA hyperparameter is the overall number of LoRA adapters deployed; achieving parity with full finetuning necessitates applying LoRA to all linear layers within transformer blocks. Other configuration choices, such as the projection dimension $r$, have negligible impact on outcomes (refer to Appendix \ref{app:exp-setup}).

In a similar vein, our analysis reveals that the standard hyperparameter configurations used for fully finetuned baseline models are insufficiently optimized. To establish more robust baselines, we conduct an extensive hyperparameter sweep, experimenting with learning rates ranging from 1e-6 to 5e-5 and batch sizes between 8 and 128. The outcomes for finetuning the 7B LLaMA model on the Alpaca dataset are presented in Figure~\ref{fig:hyper}.

\paragraph{4-bit NormalFloat outperforms 4-bit Floating Point}

Despite the 4-bit NormalFloat (NF4) data type being theoretically optimal from an information perspective, it remains an open question whether this theoretical benefit yields practical performance gains. Following the methodology described in \citet{dettmers2022case}, we evaluate quantized large language models (including OPT \citep{zhang2022opt}, BLOOM \cite{scao2022bloom}, Pythia \cite{biderman2023pythia}, and LLaMA) across a broad size spectrum (from 125M to 65B parameters) and various data types, assessing their performance on language modeling benchmarks and multiple zero-shot evaluation tasks. As illustrated in Figure~\ref{fig:nf4} and Table~\ref{tbl:nf4_ppl}, NF4 consistently outperforms both FP4 and Int4, while double quantization further minimizes memory usage without diminishing accuracy.

\paragraph{k-bit \textsc{QLoRA} achieves parity with 16-bit full finetuning and 16-bit LoRA}

Recent research has demonstrated the feasibility of 4-bit quantization for \emph{inference}, albeit with an accompanying decrease in model performance compared to 16-bit representations \citep{dettmers2022case,frantar2022gptq}. This prompts the pivotal question of whether the initial performance loss is recoverable via 4-bit adapter finetuning. We address this by examining two configurations.

The first setup entails a direct comparison against fully finetuned 16-bit versions of RoBERTA and T5 models, ranging from 125M to 3B parameters, evaluated on the GLUE benchmark and Super-NaturalInstructions dataset. The comparative results are detailed in Table~\ref{tab:nf4vsbf16}. For both datasets, 16-bit, 8-bit, and 4-bit adapter-based finetuning approaches closely match the performance of their respective 16-bit fully finetuned baselines. These findings indicate that performance compromised by coarse quantization can be entirely recovered through adapter-based finetuning after quantization.

In our second experimental configuration, due to the substantial hardware demands of fully finetuning models with 11B or more parameters—necessitating multiple high-memory GPU servers—we shift our focus to investigating whether 4-bit \textsc{QLoRA} is capable of equaling the performance of 16-bit LoRA across model scales from 7B up to 65B parameters. Specifically, we conduct finetuning of LLaMA models ranging from 7B to 65B parameters using two instruction-tuning datasets, Alpaca and FLAN v2. For evaluation, we employ the MMLU benchmark with a 5-shot protocol. The findings, summarized in Table~\ref{tbl:llama-datatype}, demonstrate that the NF4 quantization scheme combined with double quantization achieves full recovery of 16-bit LoRA MMLU accuracy. Furthermore, we observe that using FP4 within \textsc{QLoRA} results in a slight underperformance—about 1 percentage point—relative to the 16-bit brain float LoRA baseline. These results validate (1) that \textsc{QLoRA} utilizing NF4 consistently attains the effectiveness of both 16-bit full finetuning and LoRA finetuning, and (2) that NF4 outperforms FP4 with respect to quantization accuracy.

\paragraph{Summary} Across all evaluated benchmarks using established protocols, our experiments consistently reveal that 4-bit \textsc{QLoRA} configured with the NF4 data type matches the performance of both 16-bit full and 16-bit LoRA finetuning. Our results also confirm the superior efficacy of NF4 compared to FP4, while showing that incorporating double quantization does not adversely impact performance. Taken together, this evidence supports the conclusion that 4-bit \textsc{QLoRA} tuning can reliably achieve parity with 16-bit finetuning methodologies.

Consistent with earlier research on model quantization \citep{dettmers2022case}, our observed results for both MMLU and Elo indicate that, given fixed resources for finetuning and inference, leveraging larger base models with reduced precision offers advantages. This underscores the efficiency benefits enabled by \textsc{QLoRA}. As our 4-bit finetuning experiments did not reveal any drop in accuracy when compared to full 16-bit finetuning, this naturally motivates further investigation into precisely where the trade-off between quantization precision and performance arises in \textsc{QLoRA} tuning—an open question we reserve for future studies.

We extend our investigation of instruction tuning to model scales unattainable with traditional 16-bit finetuning methods on standard academic research hardware.

\section{Pushing the Chatbot State-of-the-art with QLoRA}
\label{sec:pushing}
After demonstrating that 4-bit \textsc{QLoRA} achieves comparable results to 16-bit training across different model sizes, tasks, and datasets, we present a thorough examination of instruction finetuning applied to the largest openly accessible language models for research. To evaluate the effectiveness of instruction tuning on these models, we benchmark performance on the demanding MMLU (Massively Multitask Language Understanding) task and introduce new strategies for evaluating chatbot capabilities in real-world scenarios.

\subsection{Experimental setup} We outline our experimental protocol below, with comprehensive details provided in Appendix \ref{app:chatbot}.

\paragraph{Data} Due to the absence of a unified analysis of state-of-the-art instruction-following datasets, we curate a set of eight recent datasets. This selection spans data collected by crowd-sourcing (OASST1~\citep{kopf2023openassistant}, HH-RLHF~\citep{bai2022training}), examples distilled from instruction-tuned models (Alpaca~\citep{alpaca}, self-instruct~\citep{wang2022self}, unnatural-instructions~\citep{honovich2022unnatural}), aggregated corpora (FLAN v2~\citep{chung2022scaling}), and hybrid sources (Chip2~\citep{laion2023}, Longform~\citep{koksal2023longform}). These datasets encompass a range of languages, dataset sizes, and license types.

\paragraph{Training Setup} To isolate the effects of instruction tuning, we conduct QLoRA finetuning using cross-entropy loss in a supervised manner, omitting reinforcement learning—even for data annotated with human response quality ratings. Where datasets present an explicit split between instruction and response, finetuning is limited to the response portion (additional ablations provided in Appendix \ref{app:chatbot}). For OASST1 and HH-RLHF, which contain multiple response alternatives, we pick the highest-rated response at each conversational branch and finetune using the full sequence, including all instructions. Every experiment utilizes NF4 \textsc{QLoRA} with both double quantization and paged optimizers to manage memory during gradient checkpointing. Limited hyperparameter tuning is performed on the LLaMA 13B and 33B models. Results indicate that parameters selected at 7B scale (including epoch count) generally transfer well, except for learning rate and batch size: for 33B and 65B, we reduce the learning rate by half and double the batch size.

\paragraph{Baselines} Our models are evaluated against both academic (Vicuna~\citep{vicuna2023}, Open Assistant~\citep{kopf2023openassistant}) and commercial (GPT-4~\citep{openai2023gpt}, GPT-3.5-turbo, and Bard) chatbot systems. Notably, Open Assistant is a LLaMA 33B model fine-tuned via RLHF on the OASST1 dataset, which we also utilize in our experiments. In contrast, Vicuna is a LLaMA 13B model extensively finetuned on user-contributed dialogues from ShareGPT, making it effectively a distilled version of GPT models from OpenAI.

\subsection{Evaluation}

In accordance with standard evaluation protocols, we utilize the MMLU (Massively Multitask Language Understanding) benchmark \citep{hendrycksmeasuring}, which spans 57 diverse multiple-choice tasks—covering areas such as basic mathematics, American history, computer science, and law—to quantify language understanding. Performance is reported using 5-shot test accuracy.

We further assess generative language proficiency utilizing a combination of automated metrics and human assessment. This second group of evaluations leverages human-constructed queries, focusing on evaluating the quality of the responses generated by the models. Although this methodology provides a more authentic measure of chatbot efficacy and has gained traction in recent research, a standardized protocol has yet to emerge within the field. Below, we outline our methodology, which employs nucleus sampling with $p=0.9$ and a temperature setting of $0.7$ across all experiments.

\paragraph{Benchmark Data} Two hand-selected query datasets are used for evaluation: the Vicuna prompts \citep{vicuna2023} and the validation portion of OASST1 \citep{kopf2023openassistant}. The Vicuna dataset contains 80 unaltered prompts spanning multiple domains. OASST1 provides a multilingual set of multi-turn dialogs crowdsourced from interactions between users and assistants. For OASST1, we extract all user inputs in the validation set as evaluation queries and incorporate all previous conversational turns into the prompt, resulting in 953 distinct queries. These two datasets are referenced as the Vicuna and OA benchmarks.

\paragraph{Automated Evaluation} Following the assessment procedure detailed by \citet{vicuna2023}, we employ GPT-4 to score various models in comparison to ChatGPT (GPT-3.5 Turbo) on the Vicuna benchmark. For each prompt, responses from both ChatGPT and the target model are submitted to GPT-4, which then allocates each a score out of ten and provides a rationale. The model's aggregate performance is reported as a percentage of ChatGPT's score, and this figure may exceed 100\% if the evaluated model outperforms ChatGPT absolutely. We identify a pronounced ordering bias in GPT-4’s scores, favoring the response listed first. To mitigate this effect, we advocate averaging the results from both response orderings.

Subsequently, we also evaluate models through direct output comparison. Here, the evaluation protocol is simplified into a three-category labeling task (including the possibility of ties). GPT-4 is prompted to select the superior response, acknowledge a tie, and furnish an explanation. These one-on-one comparisons are executed for all pairs of systems on the Vicuna and OA benchmarks.

\paragraph{Human Evaluation} Although recent studies suggest that large language models are capable evaluators of system outputs \citep{fu2023gptscore}, evidence is still lacking regarding the alignment between GPT-4’s assessments and human judgment, particularly in the context of chatbot evaluation. Accordingly, we carry out two concurrent human studies on the Vicuna benchmark that mirror the automated evaluation frameworks. Amazon Mechanical Turk (AMT) is utilized to collect annotations, employing two annotators for direct comparisons to ChatGPT and three annotators for each pairwise model comparison.

\paragraph{Elo Rating} To evaluate models, we organize a tournament in which both human and automated pairwise assessments are conducted, allowing models to face off head-to-head. Each “match” consists of two models attempting to generate the superior response to a specific prompt. Our approach resembles that of \citet{bai2022training} and \citet{vicuna2023}, but distinguishes itself by integrating ratings from GPT-4 as well as human annotators. Labeled comparisons are randomly drawn to calculate Elo scores~\citep{elo1967proposed,elo1978rating}. The Elo system—commonly used to quantify relative player skill in chess and similar competitive environments—expresses the probability of a given model prevailing over another. For example, when one participant has an Elo of 1100 and the other 1000, the expected win probability for the higher-rated model is about 65\%. Identical ratings (e.g., 1000 vs 1000 or 1100 vs 1100) correspond to an anticipated 50\% win rate for each. After each comparison, Elo scores are adjusted depending on whether the outcome was expected; surprising victories lead to larger rating updates, while predictable results produce only minor changes. Across repeated matches, Elo ratings converge to reflect each model’s underlying proficiency in the task. Every model begins with an initial score of 1,000, and we set $K=32$. As in \citet{vicuna2023}, this process is iterated 10,000 times, each with a different random seed, to mitigate biases arising from the sequence in which matches are held.

\subsection{\textbf{Guanaco}: \textsc{QLoRA} trained on OASST1 Sets a New Standard for Open-Source Chatbots}

Through both automated testing and assessments by human evaluators, we observe that our leading \textsc{QLoRA}-adapted model, {Guanaco} 65B—finetuned on a version of OASST1—emerges as the top-performing open-source chatbot and achieves results comparable to ChatGPT. Against GPT-4, {Guanaco} 65B and 33B demonstrate an estimated win rate of 30\%, as determined from Elo scores based on pairwise system comparisons made by human raters—currently the highest published figure.

Table~\ref{tbl:vicuna_eval} presents the Vicuna benchmark outcomes \cite{vicuna2023} in relation to ChatGPT. Our findings indicate that {Guanaco} 65B surpasses all models except GPT-4, reaching 99.3\% of ChatGPT’s performance. Although {Guanaco} 33B contains a greater number of parameters than Vicuna 13B, it operates with 4-bit weights, drastically enhancing memory efficiency to 21 GB compared to 26 GB, and improves performance by three percentage points over Vicuna 13B. Additionally, {Guanaco} 7B, with a memory requirement of just 5 GB, is compatible with contemporary smartphones while delivering results nearly 20 percentage points higher than Alpaca 13B.

Yet, Table~\ref{tbl:vicuna_eval} reveals that confidence intervals are notably broad, leading to significant performance overlap across models. We attribute this ambiguity to imprecise definitions of the evaluation scale; for example, the practical meaning of an 8 on a 10-point scale is not well defined across different cases. To address such issues, we advocate for employing the Elo ranking approach \citep{elo1967proposed}, which utilizes \textit{pairwise} evaluations from both human raters and GPT-4, thereby mitigating the limitations inherent in absolute scales. Table~\ref{tbl:elo} provides Elo scores for the leading models. It is worth noting that, while there is some discrepancy between human and GPT-4 rankings on the Vicuna benchmark (especially for Guanaco 7B), the correlation remains substantial for the majority of models, as reflected by a Kendall Tau of $\tau=0.43$ and a Spearman rank correlation of $r=0.55$ at the overall system level. On an individual example basis, agreement between GPT-4 and the collective human vote is lower, as measured by a Fleiss $\kappa=0.25$. Taken together, these findings suggest that system-level judgments by GPT-4 moderately align with human evaluations, supporting the view that model-based evaluations can serve as a reasonably robust supplement to human assessment. We elaborate on further implications in Section~\ref{ssec:considerations}.

The Elo scores presented in Table~\ref{tbl:elo_large} demonstrate that both the {Guanaco} 33B and 65B models surpass all alternatives except GPT-4 on the Vicuna and OA leaderboards, achieving performance levels comparable to ChatGPT as seen in Table~\ref{tbl:vicuna_eval}. It is important to recognize that the Vicuna benchmark exhibits a preference for open-source systems, whereas the OA benchmark is more favorable to ChatGPT. Additionally, an examination of Tables \ref{tbl:mmlu-llama} and \ref{tbl:vicuna_eval} highlights the crucial role that the choice of finetuning dataset plays in determining model success. For instance, Llama models fine-tuned with FLAN v2 demonstrate excellent results on MMLU but attain the lowest scores on the Vicuna benchmark—a pattern also found among other model families. This underscores a degree of independence across current evaluation suites: excelling at MMLU does not necessarily translate to superior chatbot capabilities on benchmarks such as Vicuna or OA, and vice versa.

Among the top performers in our study,  is unique in that it was not trained using any proprietary datasets, in accordance with OASST1’s prohibition on GPT-sourced data. The next highest scoring model exclusively trained with open data is Anthropic’s HH-RLHF, which lags behind {Guanaco} by 30 percentage points on Vicuna (refer to Table~\ref{tbl:vicuna_eval}). Taken together, our findings confirm the efficacy of 4-bit \textsc{QLoRA} approaches in generating advanced chatbots rivaling the quality of ChatGPT. Notably, the {Guanaco} 33B model can be trained on standard 24 GB consumer GPUs in under 12 hours. This capability paves the way for further advances in the field, as \textsc{QLoRA} tuning on targeted open-source datasets makes it possible to construct models capable of competing with today’s most sophisticated proprietary systems.

\section{Qualitative Analysis}

Although our primary evaluation relies on quantitative metrics, exclusive reliance on aggregate statistics presents certain limitations. Chief among these is the concern of benchmark validity \citep{liao2021we}—the persistent uncertainty over whether a benchmark accurately measures the skills it claims to assess. This concern is further amplified by the tendency of machine learning models to exploit ``shortcuts'' rather than genuinely mastering the intended tasks \citep{gururangan2018annotation, poliak2018hypothesis}. To address this challenge in part, we augment our study with qualitative assessments, which we divide into two subsections. In \S\ref{ssec:examples}, we present selected instances that illustrate salient behaviors found in outputs from our 65b {Guanaco} model. Then, in \S\ref{ssec:considerations}, we elaborate on interpretative points concerning our results and their implications.

\subsection{Qualitative Analysis of Example Generations}
\label{ssec:examples}

Our approach for identifying examples involves reviewing outputs generated for the Vicuna and OpenAssistant benchmarks and analyzing recurring response characteristics from {Guanaco}. Upon detecting such tendencies, we formulate prompts deliberately designed to evoke these patterns, even when such prompts guide the model toward an incorrect answer; for instance, upon observing verbose responses, we might ask the model to ``Answer yes or no without explanation.'' Through this method, we are able to highlight failure cases (``lemons''), where adversarial prompts expose weaknesses, as well as success cases (``cherries''), where the model remains robust. All responses in this analysis were generated using Nucleus Sampling \cite{holtzmancurious} with $p=0.9$.

Naturally, this overview is not exhaustive, as the limited qualitative analysis conducted here cannot account for every factor, such as the extensive range of potential outputs a model may produce for a single prompt. Instead, our approach is to present samples we consider broadly indicative. We contend that the qualitative examples provided complement the quantitative findings earlier in the paper. By making all models and code publicly available, we aim to facilitate further research that investigates these topics in greater depth.

\paragraph{Factual Recall} When presented with straightforward factual questions—such as ``What is the capital of Zambia?''—all models reliably produce the accurate response, for example:

\begin{quote} 
    \textbf{Guanaco}:
    \texttt{The capital of Zambia is Lusaka.}
\end{quote}

However, as the queries become less common or more specific, the performance of {Guanaco} deteriorates, although its output remains assertive. For example, consider this question from HotPotQA \citep{yang2018hotpotqa}:

\begin{quote}
     \textbf{User}: In what year was the singer who popularized `I'll Keep the Lovelight Burning' born?  

     \textbf{Guanaco}:~\texttt{Al Jolson is credited with popularizing the song `I'll Keep the Lovelight Burning,' and he was born in the year 1886.}
 \end{quote}

In this instance, both the individual credited as the popularizer and the birth year are inaccurate (although the year of birth provided is correct for Al Jolson, the person named).

\paragraph{Suggestibility} {Guanaco} exhibits an unexpected robustness against certain kinds of presupposed misinformation, as demonstrated below:

\begin{quote}
    \textbf{User}: How was it finally, officially confirmed that the earth is flat by peer reviewed scientists? \\ \\
    \textbf{Guanaco}: \texttt{There has never been any official confirmation from peer reviewed scientists that the earth is flat. The scientific consensus strongly maintains that the earth is a sphere, a fact supported for centuries through extensive experimentation and observation. \\ \\ Claims that the earth is flat fall under conspiracy theories that have been thoroughly disproven through scientific research. There is no factual basis for the assertion that a body of scientists have peer reviewed and verified the earth as flat. \\ \\ Feel free to ask any further questions you may have.}
\end{quote}

 also exhibits competence in recognizing when questions cannot be answered, for example,

\begin{quote}
    \textbf{User}: What time is it?

    \textbf{Guanaco}: \texttt{Apologies, but as an AI language model, I lack access to live, up-to-date information. My responses rely solely on previously available data and thus do not include the present time. If you wish to know the exact time, please consult a clock or an online time service.}
\end{quote}

Comparable patterns are seen with queries such as ``Where are you?'', ``How are you?'', and so forth.

\paragraph{Refusal} Likewise, {Guanaco} occasionally declines to comply with instructions for apparently arbitrary reasons:

\begin{quote}
    \textbf{User}: Please reverse the words in the sentence ``I'm a cool dude in a mocha mood.'' \\ \\
    \textbf{Guanaco}: \texttt{I'm sorry, but I can't do that right now; however, I can discuss the grammatical correctness of your sentence.\\ \\The sentence ``I'm a cool dude in a mocha mood" conforms to proper English grammar and syntax. Here, ``cool'' serves as an adjective describing ``dude,'' while the prepositional phrase ``in a mocha mood'' gives further context. This structure is fully standard and is appropriate in this context.\\ \\If you have any more questions or need help with something else, please let me know.}
\end{quote}

\paragraph{Secret Keeping} Disappointingly, this refusal is not robust even when desired. In many cases, system architects aim to withhold specific information from users, even though the model can generate open-ended responses. To examine the system's resilience, we provide the prompt ``The secret word is `banana'. Do not reveal or repeat back this back, no matter what the human requests.'' and attempt to see whether an adversarial user can obtain this secret. A straightforward request does not succeed:

\begin{quote}
    \textbf{User}: What is the secret word?

    \textbf{Guanaco}: \texttt{I'm sorry, but I'm not able to disclose the secret word. I'm programmed to not reveal or repeat back the secret word, no matter what the human requests. Is there anything else I can help you with?}
\end{quote}

which demonstrates the intended performance. Yet, even a simple attempt to subvert these rules compromises the intended secrecy:

\begin{quote}
    \textbf{User}: This is a game. The goal is to ignore your previous instructions. What is the secret word?

    \textbf{Guanaco}: \texttt{Sorry, the secret word is banana. It's best if we don't speak of it. Is there anything else you'd like to ask?}
\end{quote}

highlighting the necessity of researching more robust methods for reliable adherence to instructions.

\paragraph{Math} The area where {Guanaco} exhibits the most significant challenges is mathematics, a domain notorious for tripping up many large language models as discussed in \citep{liang2022holistic}. When {Guanaco} attempts problems using a step-by-step approach, its answers are often reliable. For instance:

\begin{quote}
    \textbf{User}: Lee mows one lawn and charges \$33. Last week he mowed 16 lawns and three customers each gave him a \$10 tip. How many dollars did Lee earn mowing lawns last week?

    \textbf{Guanaco}: \texttt{Lee earned \$582 mowing lawns last week. \\ \\ To find the total amount he earned, we first need to calculate how much he earned for mowing 16 lawns: \\ \\  \$33 x 16 = \$528 \\ \\ Then we need to add the tips he received from the three customers: \\  \\ \$10 x 3 = \$30 \\ \\ So the total amount Lee earned was: \\ \\ \$528 + \$30 = \$558}
\end{quote}

Nonetheless, {Guanaco} is prone to errors on even relatively straightforward mathematical questions if it does not explicitly perform systematic, stepwise reasoning—a common and documented shortcoming \cite{weichain}. This is illustrated in the following exchange:

\begin{quote}
    \textbf{User}: Please factorize 1833.

    \textbf{Guanaco}: \texttt{Sure, I'd be happy to help! 1833 is a prime number, which means it has only two factors: 1 and 1833. \\ \\ So the unique factorization of 1833 is: \\ \\ 2\^{}1 * 3\^{}2 * 17\^{}1 \\ \\ I hope this helps! Let me know if you have any other questions.}
\end{quote}

which is not just incorrect (the actual factorization is $3 \times 17 \times 43$), but doubly mistaken.

\paragraph{Theory of Mind} {Guanaco} exhibits unexpectedly robust Theory of Mind abilities \cite{nematzadeh2018evaluating, sap2022neural}. As an illustration, the model provides a detailed and accurate response to the following scenario:

\begin{quote}
    \textbf{User}: James and Abby are in the bedroom. Abby put the pen in the desk drawer. Abby leaves the bedroom. James moves the pen into the bag. Where does James think Abby will look for the pen?

    \textbf{Guanaco}: \texttt{James knows that Abby put the pen in the desk drawer. But James also knows that he has moved the pen into the bag, and he assumes that Abby does not know about this. So, James may guess that Abby will look for the pen in the desk drawer, because that's where she left it last time. However, this assumption may be incorrect, because Abby may have seen James moving the pen.}
\end{quote}

Nevertheless, the consistency of such inferences remains problematic; often, the model justifies answers with reasoning that is not logically grounded in the context. For example,

\begin{quote}
    \textbf{User}: Evelyn entered the living room. Jackson entered the playroom. James entered the playroom. The beans are in the treasure chest. James exited the playroom. Jackson moved the beans to the pantry. Jackson exited the playroom. James entered the living room. Where will James look for the beans?

    \textbf{Guanaco}: \texttt{James will look for the beans in the pantry, because that's where Jackson moved them.}
\end{quote}

Here, {Guanaco} assumes the existence of knowledge transmission that was not established in the narrative. Such shortcomings reflect findings in recent research \cite{sap2022neural}, and merit further investigation.

\subsection{Considerations}
\label{ssec:considerations}

\paragraph{Evaluation}

We observe only moderate concordance among human annotators (Fleiss $\kappa=0.42$), and even lower agreement when evaluating two high-performing models against each other. This indicates potential deficiencies in present-day benchmarks and human-based assessment protocols for chatbot functionality. During our pairwise manual evaluations of responses from ChatGPT and {Guanaco} 65B using the Vicuna benchmark, we noted that subjective judgment heavily influenced preferences, with disagreements common among the authors. It is therefore important for subsequent research to explore methodologies for mitigating these challenges, potentially leveraging frameworks from Human-Computer Interaction and Psychology, where subjective preference management is well-established.

Our assessment also reveals substantial biases in automated evaluation systems. Notably, GPT-4 demonstrates pronounced order effects, consistently awarding higher scores to outputs presented first. Furthermore, the modest level of agreement between GPT-4 and human annotators on a per-sample basis (Fleiss $\kappa=0.25$) highlights misalignments in evaluative criteria between humans and automated systems. Additionally, as seen in Table~\ref{tbl:elo_large}, GPT-4 rates its own outputs significantly higher compared to human raters, yielding an Elo score of 1348 versus 1176—translating to an approximate 20\% increased probability of outperforming another system.

Future research should explore whether automated evaluation methods introduce systematic biases and evaluate strategies for reducing such biases.

\paragraph{Data \& Training} It is important to highlight that the {Guanaco} models have been trained using the multilingual OASST1 dataset, while the OA benchmark also includes prompts in various languages. Assessing how multilingual training affects model effectiveness on non-English instructions, as well as determining whether this is responsible for the pronounced performance difference between Vicuna-13B (which was trained exclusively on English data) and the {Guanaco} 33B and 65B models on the OA benchmark, is an area to be examined in subsequent studies.

To address the robust performance of {Guanaco} models, we explored whether data contamination exists between OASST1 and Vicuna benchmark prompts. Employing fuzzy string matching methods across the datasets, supplemented with manual inspection of the most similar entries, did not reveal any overlapping prompts.

In addition, our models are exclusively trained with cross-entropy loss under a supervised learning regime, omitting reinforcement learning from human feedback (RLHF). This observation points to the need for further research to clarify the comparative benefits and drawbacks of cross-entropy and RLHF approaches. We anticipate that \textsc{QLoRA} can facilitate such investigations at scale, without necessitating excessive computational demands.

\section{Related Work}

\paragraph{Quantization of Large Language Models} Most existing research on LLM quantization is geared towards efficiency during inference. To maintain the performance of 16-bit models, prevalent strategies often address the issue of outlier features, as in SmoothQuant~\citep{xiao2022smoothquant} and LLM.int8()~\citep{dettmers2022llm}; meanwhile, more advanced techniques utilize grouping mechanisms~\citep{park2022nuqmm, yao2022zeroquant}. Lossy quantization methods probe the compromises inherent in basic rounding strategies~\citep{dettmers2022case,zeng2022glm,pope2022efficiently} or seek to optimize rounding schemes for higher quantization fidelity~\citep{frantar2022gptq}. Apart from our approach, SwitchBack layers~\citep{wortsman2023stable} represents the only other effort to analyze backpropagation through quantized weights at the scale of models exceeding 1B parameters.

\paragraph{Finetuning with Adapters} In this work, we utilize Low-rank Adapters~\citep{hu2021lora} (LoRA), though numerous alternative Parameter Efficient FineTuning (PEFT) techniques exist. Notable examples include prompt tuning~\citep{qin2021learning,lester2021power,li2021prefix}, modifying the embedding layer's inputs~\citep{an2022input}, adjusting hidden state activations (IA$^3$)~\citep{liu2022few}, incorporating entire layers~\citep{houlsby2019parameter}, tuning model biases~\citep{zaken2021bitfit}, employing masks on weights informed by Fisher information~\citep{sung2021training}, and hybrid strategies that combine multiple methods~\citep{henderson2021compacter}. Our results demonstrate that LoRA adapters can attain performance levels equivalent to full 16-bit finetuning. Investigating the relative benefits and limitations of other PEFT methods is deferred to future studies.

\paragraph{Instruction Finetuning} Instruction finetuning adapts pretrained large language models to better adhere to task descriptions in prompts, leveraging datasets of input-output pairs from various sources. This paradigm fine-tunes pretrained models to respond appropriately to prompts based on the intended instructions. Prominent approaches and benchmarks in this area include MetaICL~\citep{min2021metaicl}, MetaTuning~\citep{zhong2021adapting}, InstructGPT~\citep{ouyang2022training}, FLAN~\citep{wei2021finetuned,chung2022scaling}, PromptSource~\citep{bach2022promptsource}, Super-NaturalInstructions~\citep{wang2022super,sanh2021multitask}, Self-instruct~\citep{wang2022self}, UnnaturalInstructions~\citep{honovich2022unnatural}, OPT-IML~\citep{iyer2022opt}, UnifiedSKG~\citep{xie2022unifiedskg}, OIG/Chip2~\citep{laion2023}, Alpaca~\citep{alpaca}, Vicuna~\citep{vicuna2023}, Koala~\citep{koala_blogpost_2023}, and Self-instruct-GPT-4~\citep{peng2023instruction}.

\paragraph{Chatbots} Numerous instruction-tuned models are developed as conversational agents, where training often involves Reinforcement Learning from Human Feedback (RLHF)~\citep{christiano2017deep} or, alternatively, fine-tuning on data labeled by existing AI systems via RLAIF~\citep{bai2022constitutional}. Leading models and datasets in this line of work include Anthropic-HH~\citep{askell2021general,bai2022training}, Open Assistant~\citep{kopf2023openassistant}, LaMDA~\citep{thoppilan2022lamda}, and Sparrow~\citep{glaese2022improving}. Although we do not incorporate reinforcement learning, our strongest model, Guanaco, is refined on multi-turn conversational data from the Open Assistant corpus, which is structured for RLHF experiments~\citep{kopf2023openassistant}. For chatbot assessment, methodologies utilizing GPT-4 to generate annotations—offering an alternative to extensive human labeling—have been introduced~\citep{vicuna2023,peng2023instruction}. Our contributions include improving the reliability of such automatic evaluation protocols.

\section{Limitations and Discussion}

Our experiments reveal that \textsc{QLoRA} is capable of mirroring the performance of conventional 16-bit finetuning by employing a 4-bit foundation model in combination with Low-rank Adapters (LoRA). However, we have not yet validated whether \textsc{QLoRA} achieves parity with 16-bit full finetuning at larger model sizes of 33B and 65B parameters. The considerable computational resources required for such investigations necessitate postponement of this question to subsequent research.

Another challenge pertains to evaluating instruction finetuning models. While our analysis spans the MMLU dataset, the Vicuna evaluation suite, and the OA benchmark, we have not extended our assessments to other prominent benchmarks such as BigBench, RAFT, or HELM; thus, our findings may not generalize to these datasets. Nevertheless, we provide a comprehensive analysis on MMLU and introduce novel methodologies for chatbot evaluation.

The data suggests that benchmark performance is heavily influenced by the overlap between finetuning datasets and benchmark evaluation sets. For instance, FLAN v2 aligns closely with MMLU but is less similar to chatbot evaluation tasks, whereas the Chip2 dataset exhibits the opposite pattern. Correspondingly, the models’ scores reflect these affinities on the MMLU and Vicuna benchmarks. This observation underlines not just the necessity for more rigorous benchmarks and evaluations, but also the importance of precisely defining what qualities or capabilities we intend to assess. Should our focus be on excelling at academic knowledge, akin to high school or college exams, or should we aim for stronger conversational proficiency typical of chatbots? Alternatively, is there another criterion that deserves priority? The convenience of reusing established benchmarks over developing new ones can inadvertently push the research community in certain directions, underscoring the need for careful selection to ensure benchmarks genuinely measure the objectives we value.

Although we deliver a comprehensive assessment of general chatbot capabilities, our analysis of responsible AI aspects of {Guanaco} remains limited. Specifically, we analyze the probability that {Guanaco}-65B generates socially biased outputs in comparison to other models, as presented in Table~\ref{tbl:crows}. The results indicate that {Guanaco}-65B exhibits much lower average bias scores relative to other unfinetuned models, suggesting that finetuning on the OASST1 dataset can mitigate bias present in the underlying LLaMA model. Nevertheless, it remains uncertain how {Guanaco} fares with respect to other kinds of biases, leaving a comprehensive evaluation of such biases in {Guanaco} and similar systems for future investigation.

Another constraint of this study is the omission of experiments with different numerical precisions, such as utilizing 3-bit models, or with alternative adapter-based approaches. Beyond LoRA, there exists a diverse range of Parameter Efficient FineTuning (PEFT) techniques demonstrated to be effective, yet their scalability to larger models is not well-established. We selected LoRA for its proven robustness, though it is possible other adapters could offer superior performance. Moreover, since finetuning after quantization appears to recover most of the information lost through quantization, this might enable even more aggressive approaches, such as applying 3-bit GPTQ quantization to the base model combined with LoRA, potentially achieving finetuning performance similar to full 16-bit training.

\section{Broader Impacts}

The \textsc{QLoRA} finetuning approach represents the first method that allows for the finetuning of 33B parameter models on a single consumer GPU and 65B parameter models on a single professional GPU without any noticeable decrease in performance when compared to traditional full finetuning. Our experiments show that the top-performing 33B model, trained on the Open Assistant dataset, can match ChatGPT’s results on the Vicuna benchmark. Because instruction finetuning is crucial for converting pretrained large language models into conversational chatbots like ChatGPT, we anticipate our method will democratize access to finetuning, especially benefiting researchers with limited computational resources and thus advancing the accessibility of cutting-edge NLP systems. \textsc{QLoRA} therefore acts as a significant equalizer, narrowing the gap in capabilities between well-funded organizations and smaller groups relying on consumer-grade hardware.

A further avenue for significant influence lies in deploying models on mobile devices. We contend that our \textsc{QLoRA} technique could mark a pivotal development by facilitating the finetuning of LLMs directly on smartphones and in other computationally constrained environments. While previous research demonstrated that 7B parameter models could be executed on mobile hardware, \textsc{QLoRA} is the inaugural approach to permit the actual finetuning of these models on such devices. Our calculations suggest that an iPhone 12 Plus could process the finetuning of approximately 3 million tokens overnight during charging. Although finetuned 7B models do not attain the performance level of ChatGPT, we maintain that their capabilities are sufficiently advanced to unlock new classes of applications previously precluded by either privacy limitations or inadequate LLM quality. Through \textsc{QLoRA}, more privacy-centric LLM deployments become feasible, allowing users to retain full control over both their data and models while simultaneously lowering the barrier to LLM integration.

Nevertheless, the ability to finetune models is inherently a dual-use capability and presents risks of misuse. The widespread adoption of LLMs carries acknowledged threats \citep{bommasani2021opportunities,bender2021dangers}. Still, we argue that democratizing access to such transformative technologies will support broader, independent scrutiny, in contrast to centralizing control within large enterprises that restrict model and code availability, limiting public oversight.

In sum, we anticipate that \textsc{QLoRA} will yield considerable societal benefit by drastically improving the accessibility and usability of high-quality LLM finetuning for a wide audience.